% problem-000247 8 光化学反应与重排
\section*{8. 光化学反应与重排}
光激发能让有机分子发生特殊的化学转化。例如，羰基可被光激发形成双自由基活性中间体，进而诱发碳碳键的断裂。下图所示是最常见的两种转化模式，被称为Norrish-1 型和 Norrish-2 型反应。

（图：）

\subsection*{8-1}
碳氮双键也可被光激发，发生光诱导化学转化。画出下述反应的关键中间体。

（图：）

\subsection*{8-2}
研究发现，一些A的类似化合物在非光照的常规条件下也能发生重排反应。如下图所示，化合物C和D仅有一个甲基取代的区别，但在甲醇解反应中却分别得到两种不同的产物E和F。

（图：）

D

\subsection*{8-2}
-1 画出 C 到E经历的关键中间体。

\subsection*{8-2}
-2画出D到F经历的关键中间体。

\subsection*{8-3}
哌啶化合物G在光照下可转化为环戊胺H，画出下述反应的关键中间体。

（图：）

第38届中国化学奥林匹克（决赛）理论第一场试题 2024年10月26日广州

\subsection*{8-4}
苯甲酰硅化合物J在光照下并未发生Norrish-1 型断裂，而是生成一种与J化学式相同且含有六电子碳原子的活性中间体，画出该中间体结构。

（图：）

\subsection*{8-5}
下列光反应可生成α-甲基苯乙烯、苯甲醛及化合物K，高分辨质谱显示K的化学式为 $C _ { 1 6 } \mathrm { H } _ { 1 6 } \mathrm { O }$ 。画出K的结构，并注明立体化学。

（图：）
